\chapter*{Abstract}
\addcontentsline{toc}{chapter}{Abstract}

Software-defined networking (SDN) controllers can use a network model to estimate per-flow delay before changing routes. Graph neural network (GNN) models such as RouteNet-Fermi are accurate but costly, and depend on how well their training data cover the networks met in operation. This project investigates whether a smaller GNN, trained with help from a larger one, can give usable delay estimates at lower cost. Training RouteNet-Fermi on 10- to 170-node instead of only 10-node topologies lowered the mean absolute percentage error (MAPE) on 478 unseen 200- to 300-node scenarios from 8.40\% to 5.37\%. A lightweight student, Queue-Link Fusion, merges each link with its queues, cutting message passing from 24 to 8 stages and parameters from 27,025 to 4,505. Its forward pass was 4.4 to 7.9 times faster on a CPU, but end to end the saving was only about 8\%, because building the input graph took about 90\% of the time. The student's error was about twice the teacher's. A first distillation run lowered it from 11.92\% to 9.83\%, but a controlled re-test with three paired seeds, identical initial weights and a relative distillation loss showed no consistent benefit: the distilled student was better in two pairs and worse in one on validation, and \pending{test result pending}. The student's training error (7.1\%) already exceeded the teacher's validation error (4.7\%), so its limit is capacity rather than guidance. In a six-switch Mininet emulation with a Ryu controller, all three frozen models ranked three candidate paths in the measured order of delay, and the chosen path had 8.8\% lower delay than the hop-count default. An Abilene emulation failed because its traffic lay far outside the training range. Topology diversity improved accuracy, but compression cost more than diversity recovered, and graph construction, not the model, dominated the cost.

\vspace{4mm}
\noindent\textbf{Keywords:} software-defined networking, graph neural networks, network modelling, knowledge distillation, Mininet, out-of-distribution generalisation
